% ============================================================================
% SUPPLEMENTARY MATERIALS: POLYPHARMACY RISK ASSESSMENT
% ============================================================================

\documentclass[11pt,a4paper]{article}
\usepackage[utf8]{inputenc}
\usepackage{amsmath,amssymb}
\usepackage{booktabs}
\usepackage{geometry}
\usepackage{graphicx}
\usepackage{hyperref}
\usepackage{float}
\usepackage{longtable}
\geometry{margin=1in}

\begin{document}

\title{Supplementary Materials: Polypharmacy Risk Assessment}
\date{}
\maketitle

\tableofcontents
\newpage

\section{Polypharmacy Risk Index (PRI) Derivation}

\subsection{Network-Based Risk Metrics}

The Polypharmacy Risk Index integrates four complementary network centrality measures that capture different aspects of drug interaction risk:

\subsubsection{Degree Centrality}
\begin{equation}
C_{\text{degree}}(d) = \frac{\text{deg}(d)}{N - 1}
\end{equation}
where $\text{deg}(d)$ is the number of edges connected to drug $d$ and $N$ is the total number of drugs in the network.

\subsubsection{Weighted Degree Centrality}
\begin{equation}
C_{\text{weighted}}(d) = \frac{\sum_{j \in \Gamma(d)} w_{dj}}{\sum_{i,j} w_{ij}}
\end{equation}
where $\Gamma(d)$ is the set of neighbors of $d$ and $w_{dj}$ is the severity weight of the interaction between $d$ and $j$.

\subsubsection{Betweenness Centrality}
\begin{equation}
C_{\text{betweenness}}(d) = \sum_{s \neq d \neq t} \frac{\sigma_{st}(d)}{\sigma_{st}}
\end{equation}
where $\sigma_{st}$ is the number of shortest paths between $s$ and $t$, and $\sigma_{st}(d)$ is the number of those paths passing through $d$.

\subsubsection{Severity Profile Score}
\begin{equation}
S(d) = \frac{|\{j : (d,j) \in E \land \text{severity}(d,j) \in \{\text{Contra}, \text{Major}\}\}|}{|\{j : (d,j) \in E\}|}
\end{equation}

\subsection{Weight Optimization}

The PRI weights (0.25, 0.30, 0.20, 0.25) were optimized through grid search to maximize correlation with clinical severity assessments while maintaining balanced contribution from each metric.

\section{Risk Escalation Analysis}

\subsection{Polypharmacy Thresholds}

\begin{table}[H]
\centering
\caption{Risk Escalation by Number of Concurrent Drugs}
\label{tab:risk_escalation}
\begin{tabular}{lcccc}
\toprule
\textbf{Drug Count} & \textbf{Avg. Interactions} & \textbf{Severe \%} & \textbf{Mean PRI} & \textbf{Risk Level} \\
\midrule
2 & 0.8 & 5.2\% & 0.12 & Low \\
3 & 2.1 & 8.4\% & 0.21 & Low-Medium \\
4 & 4.3 & 12.7\% & 0.32 & Medium \\
5 & 7.2 & 18.3\% & 0.45 & Medium-High \\
6+ & 12.1 & 26.8\% & 0.58 & High \\
\bottomrule
\end{tabular}
\end{table}

\section{High-Risk Drug Categories}

\subsection{Top 10 Highest PRI Drugs}

\begin{table}[H]
\centering
\caption{Drugs with Highest Average PRI Scores}
\label{tab:high_pri_drugs}
\begin{tabular}{llcc}
\toprule
\textbf{Drug} & \textbf{Category} & \textbf{PRI Score} & \textbf{Interaction Count} \\
\midrule
Warfarin & Anticoagulant & 0.892 & 385 \\
Phenytoin & Anticonvulsant & 0.867 & 312 \\
Rifampin & Antibiotic & 0.845 & 289 \\
Carbamazepine & Anticonvulsant & 0.823 & 275 \\
Amiodarone & Antiarrhythmic & 0.812 & 268 \\
Ketoconazole & Antifungal & 0.798 & 254 \\
Cyclosporine & Immunosuppressant & 0.785 & 241 \\
Digoxin & Cardiac Glycoside & 0.772 & 228 \\
Lithium & Mood Stabilizer & 0.758 & 215 \\
Methotrexate & Antimetabolite & 0.745 & 203 \\
\bottomrule
\end{tabular}
\end{table}

\section{Validation Against Clinical Guidelines}

The PRI thresholds were validated against published clinical guidelines for polypharmacy management:

\begin{itemize}
    \item AGS Beers Criteria\textsuperscript{\textregistered} 2019 for potentially inappropriate medications in older adults
    \item STOPP/START Criteria version 2 for medication review
    \item NHS Scotland Polypharmacy Guidance 2018
\end{itemize}

\section{Statistical Analysis}

\subsection{PRI Distribution Characteristics}

\begin{table}[H]
\centering
\caption{PRI Score Distribution Statistics}
\label{tab:pri_stats}
\begin{tabular}{lr}
\toprule
\textbf{Statistic} & \textbf{Value} \\
\midrule
Mean & 0.342 \\
Median & 0.298 \\
Standard Deviation & 0.186 \\
Interquartile Range & 0.245 \\
Skewness & 0.824 (right-skewed) \\
Kurtosis & 2.91 \\
\bottomrule
\end{tabular}
\end{table}

\bibliographystyle{plain}
\begin{thebibliography}{3}
\bibitem{beers2019} American Geriatrics Society 2019 Beers Criteria Update Expert Panel. American Geriatrics Society 2019 Updated AGS Beers Criteria. \textit{J Am Geriatr Soc}. 2019;67(4):674-694.
\bibitem{stopp2015} O'Mahony D, et al. STOPP/START criteria for potentially inappropriate prescribing in older people: version 2. \textit{Age Ageing}. 2015;44(2):213-218.
\bibitem{nhs2018} NHS Scotland. Polypharmacy Guidance. 3rd ed. Edinburgh: Scottish Government; 2018.
\end{thebibliography}

\end{document}
